\begin{mistake}{2026-05-27}{08}

\begin{problemcontent}
设 \(f(x) = |(x-1)(x-2)^2(x-3)^3|\)，则 \(f'(x)\) 不存在的点个数是（ \quad ）
(A) 0 
(B) 1 
(C) 2 
(D) 3
\end{problemcontent}

\begin{solutioncontent}
正确答案：(B)

对于含有绝对值的函数 \(f(x) = |g(x)|\)，若 \(g(x)\) 在 \(x=a\) 处可导且 \(g(a)=0\)：
1. 若 \(x=a\) 是 \(g(x)\) 的一重零点（即 \(g'(a) \neq 0\)），则 \(f(x)\) 在 \(x=a\) 处不可导；
2. 若 \(x=a\) 是 \(g(x)\) 的多重零点（即 \(g'(a) = 0\)），则 \(f(x)\) 在 \(x=a\) 处可导，且 \(f'(a) = 0\)。

令 \(g(x) = (x-1)(x-2)^2(x-3)^3\)，其零点为 \(x_1=1, x_2=2, x_3=3\)。
- \(x=1\) 的因式指数为 1（一重零点），故 \(f'(1)\) 不存在。
- \(x=2\) 的因式指数为 2（二重零点），故 \(f'(2)=0\)。
- \(x=3\) 的因式指数为 3（三重零点），故 \(f'(3)=0\)。

（严谨验证 \(x=1\) 处：\(\lim_{x \to 1^\pm} \frac{f(x)-f(1)}{x-1} = \lim_{x \to 1^\pm} \frac{\pm(x-1)(x-2)^2|x-3|^3}{x-1} = \pm 8\)，左右导数不等，故不可导。）

综上，仅在 \(x=1\) 处导数不存在，个数为 1，选 (B)。
\end{solutioncontent}

\mistakeknowledge{
\begin{itemize}
  \item \textbf{核心判定}：考察 \(f(x) = |(x-x_1)^{p_1}(x-x_2)^{p_2}\cdots|\) 的可导性，只需看各因式指数 \(p_i\)。
  \item \textbf{秒杀技巧}：若 \(p_i = 1\)，在该零点处不可导；若 \(p_i > 1\)，在该零点处可导且导数为 0。
  \item \textbf{易错点}：误认为绝对值函数在所有零点处均不可导。只有图象“穿过”x轴（一重根）折叠后才产生尖点；“相切”于x轴（多重根）折叠后依旧平滑。
\end{itemize}}

\end{mistake}
